% Unabhängigkeit – bank/unabhaengigkeit/ (zusammenbau v0.4, Heft)
\einheitenkopf[t4]{Unabhängigkeit}
\setcounter{aufgabe}{20}

% Anlauf (ohne Prüfkennung)
\begin{aufgabe}{Unabhängigkeit prüfen – Anlauf}
\begin{teile}
\teil Eine Vierfeldertafel mit allen Feldern und Rändern liegt vor. Wo stehen die drei Zahlen der Produktregel? \\ \kreuz{alle drei stehen in der Tafel: zwei Ränder und ein Feld} \\ \kreuz{nur die Ränder stehen da, der Schnitt muss berechnet werden} \\ \kreuz{keine steht da, alle drei sind zu beschaffen}
\teil Die Tafel zeigt 100 Schüler nach den Ereignissen M (Mädchen) und F (kommt mit dem Fahrrad). Prüfe mit der Produktregel, ob M und F stochastisch unabhängig sind.

\vierfeldertafel{M}{F}{20,30,50,20,30,50,40,60,100}
\teil In einem Sportverein sind 60 Mitglieder, 36 davon Erwachsene. 20 Mitglieder spielen Volleyball, darunter 8 Jugendliche. Bestimme die fehlenden Anzahlen und prüfe, ob E (Erwachsener) und V (spielt Volleyball) stochastisch unabhängig sind.
\end{teile}
\end{aufgabe}

\begin{aufgabe}{Unabhängigkeit prüfen – Prüfungsaufgaben}
\begin{teile}
\teil Die Tafel zeigt 250 Haushalte nach den Ereignissen S (Stadtwohnung) und H (hat ein Haustier). Sind S und H unabhängig? Prüfe mit der Produktregel, ob S und H stochastisch unabhängig sind. \hfill (Abitur 2024 GK)

\vierfeldertafel{S}{H}{70,105,175,30,45,75,100,150,250}
\teil In einer Region traten 12\,400 Personen zur Fahrprüfung an, davon waren 3\,100 mindestens 30 Jahre alt. 9\,920 Prüflinge bestanden, darunter 7\,750 jüngere. Prüfe, ob der Anteil der Bestandenen unter den Älteren mit dem Gesamtanteil übereinstimmt, ob Ä (mindestens 30 Jahre) und S (bestanden) stochastisch unabhängig sind, und deute das Ergebnis \hfill (Abitur 2019 GK).
\teil Ein Paketdienst befördert 600 Pakete, 150 davon sind schwer. 90 der schweren Pakete gehen ins Ausland, insgesamt gehen 240 Pakete ins Ausland. Untersuche, ob der Anteil der Auslandspakete unter den schweren Paketen ebenso groß ist wie unter den nicht schweren \hfill (Abitur 2022 GK).
\end{teile}
\end{aufgabe}

% Anlauf (ohne Prüfkennung)
\begin{aufgabe}{Als Argument – Anlauf}
\begin{teile}
\teil Eine Münze zeigt dreimal hintereinander Wappen. Wie wahrscheinlich ist beim nächsten Wurf Zahl? \\ \kreuz{$0{,}5$} \\ \kreuz{größer als $0{,}5$} \\ \kreuz{kleiner als $0{,}5$}
\teil Ein Glücksrad zeigt bei jeder Drehung mit $50\,\%$ Wahrscheinlichkeit Grün. Nach 40 Drehungen mit 27-mal Grün sagt Ben: „Jetzt kommt eine Weile seltener Grün, damit es wieder passt.“ Beurteile die Aussage.
\end{teile}
\end{aufgabe}

\begin{aufgabe}{Als Argument – Prüfungsaufgaben}
\begin{teile}
\teil Bei einem Online-Quiz beantwortet Jan jeden Samstag 8 Fragen; jede Frage beantwortet er unabhängig von den anderen mit der Wahrscheinlichkeit $0{,}6$ richtig. Die Anzahl der richtigen Antworten an einem Samstag ist binomialverteilt. An drei Samstagen hintereinander hatte er jeweils 7 richtige. Jan meint, die Chance auf 7 richtige sei an diesem Samstag deshalb deutlich kleiner. Beurteile die Aussage \hfill (Abitur 2021 GK).
\teil Eine Torwartin hält bei jedem Elfmeter unabhängig von den anderen mit der Wahrscheinlichkeit $0{,}3$. Die Anzahl der gehaltenen Elfmeter in einer Serie von fünf ist binomialverteilt. In den letzten drei Serien hielt sie jeweils keinen einzigen. Ihr Trainer sagt, in der nächsten Serie werde sie mit höherer Wahrscheinlichkeit als sonst mindestens einen halten. Beurteile die Aussage \hfill (Abitur 2021 GK).
\end{teile}
\end{aufgabe}
